\documentclass[11pt,a4paper]{article}

\usepackage[utf8]{inputenc}
\usepackage[T1]{fontenc}
\usepackage[spanish,es-tabla]{babel}
\usepackage[a4paper,margin=2.5cm]{geometry}
\usepackage{booktabs}
\usepackage{longtable}
\usepackage{array}
\usepackage{xcolor}
\usepackage{enumitem}
\usepackage{hyperref}
\usepackage{fancyhdr}
\usepackage{titlesec}

\definecolor{riesgoalto}{RGB}{178,34,34}
\definecolor{riesgomedio}{RGB}{184,134,11}
\definecolor{riesgobajo}{RGB}{34,120,60}

\hypersetup{
    colorlinks=true,
    linkcolor=blue!50!black,
    urlcolor=blue!50!black,
    pdftitle={Alternativas de respaldo ante caida de OpenAI},
    pdfauthor={Escaneria - Migasa}
}

\titleformat{\section}{\normalfont\Large\bfseries}{\thesection}{1em}{}
\titleformat{\subsection}{\normalfont\large\bfseries}{\thesubsection}{1em}{}

\pagestyle{fancy}
\fancyhf{}
\fancyhead[L]{\small Alternativas de respaldo ante caída de OpenAI}
\fancyhead[R]{\small Proyecto Escaneria (Migasa)}
\fancyfoot[C]{\thepage}

\newcommand{\riesgo}[1]{\textcolor{riesgoalto}{\textbf{[#1]}}}
\newcommand{\advertencia}[1]{\textcolor{riesgomedio}{\textbf{[#1]}}}
\newcommand{\ok}[1]{\textcolor{riesgobajo}{\textbf{[#1]}}}

\title{\textbf{Alternativas de respaldo ante una caída de OpenAI}\\[0.3em]
\large Evaluación de Azure OpenAI, OpenRouter y LiteLLM para el pipeline\\
de extracción automática de facturas (Escaneria) --- Migasa}
\author{}
\date{11 de agosto de 2026}

\begin{document}

\maketitle
\thispagestyle{empty}

\begin{abstract}
\noindent
Este informe complementa la evaluación previa de Mammouth.ai. El objetivo real
que motiva la búsqueda de una capa de acceso a varios modelos de IA no es evitar
integraciones directas por sí mismas, sino \textbf{poder seguir procesando
facturas si un día la API de OpenAI deja de estar disponible}. Bajo esa premisa,
Mammouth.ai resultó poco adecuado (riesgos de compatibilidad no confirmados,
modelo de cuotas pensado para chat interactivo, coste peor que el actual). Aquí
se evalúan tres alternativas más orientadas a uso programático de respaldo:
\textbf{Azure OpenAI}, \textbf{OpenRouter} y \textbf{LiteLLM} (autoalojado).
\end{abstract}

\tableofcontents

\section{Objetivo y criterio de evaluación}

El criterio ya no es "evitar integrarse con cada proveedor" en abstracto, sino
tres preguntas concretas para cada alternativa:

\begin{enumerate}
    \item Si hoy \texttt{extract\_invoice\_with\_agent} usa
    \texttt{response\_format} con JSON Schema en modo \texttt{strict}, y
    \texttt{imagenes.py} manda imágenes como entrada para visión,
    \textbf{¿sigue funcionando igual con esta alternativa?}
    \item Si OpenAI se cae, \textbf{¿cuánto esfuerzo (manual o automático) hace
    falta para seguir procesando facturas?}
    \item \textbf{¿El coste y el modelo de facturación son predecibles} (pago por
    consumo real, sin cuotas de sesión ni degradación silenciosa de modelo)?
\end{enumerate}

\section{Opción 1: Azure OpenAI}

\subsection{Qué es}
El mismo servicio de OpenAI (incluido \texttt{gpt-4.1}), operado por Microsoft en
su propia infraestructura de nube, con SLA formal de disponibilidad.

\subsection{Evaluación}

\begin{itemize}
    \item \ok{Sin riesgo} \textbf{Compatibilidad total.} Es la misma API que
    usáis hoy: \texttt{response\_format} con \texttt{strict} y visión funcionan
    exactamente igual, sin necesidad de probar nada nuevo en la lógica de
    extracción.
    \item \ok{Esfuerzo mínimo} El propio diseño de \texttt{OPENAI\_BASE\_URL} en
    \texttt{build\_client()} (\texttt{logic.py}) ya está pensado para apuntar a
    un endpoint alternativo como este. El cambio de proveedor hoy sería manual
    (editar \texttt{.env} y reiniciar el servicio), salvo que se añada
    detección automática de fallo.
    \item \advertencia{Cobertura parcial} Protege frente a una caída del
    \emph{endpoint} de \texttt{openai.com} (el tipo de incidencia más habitual),
    pero no frente a un problema propio del modelo GPT en sí, al ser la misma
    familia de modelo servida por otra infraestructura.
    \item \ok{Facturación} Pago por consumo real, sin cuotas de sesión ni
    degradación de modelo; tarifa por token publicada y comparable a la de
    OpenAI directo.
\end{itemize}

\section{Opción 2: OpenRouter}

\subsection{Qué es}
Agregador de modelos orientado a desarrolladores (no a chat de consumo), con más
de 300 modelos de distintos proveedores bajo una única API compatible con el
formato de OpenAI.

\subsection{Evaluación}

\begin{itemize}
    \item \ok{Facturación transparente y predecible.} Cobra el precio base del
    proveedor más una comisión fija y publicada del 5,5\% por petición. No hay
    suscripción con crédito incluido ni degradación automática de modelo: si
    hay saldo, la petición se sirve; si no, se rechaza con un error explícito.
    Esto coincide con el modelo de "pago por consumo" al que ya estáis
    acostumbrados con OpenAI directo.
    \item \ok{Selección de modelo explícita.} Se elige el modelo exacto
    (\texttt{openai/gpt-4.1}, \texttt{anthropic/claude-3.5-sonnet}...), no un
    alias que Mammouth pueda cambiar sin avisar.
    \item \ok{Fallback automático entre proveedores, como función nativa.} Se
    define una lista de modelos por prioridad; si el primero falla (error 5xx,
    rate limit, caída), OpenRouter reintenta automáticamente con el siguiente.
    Es exactamente la función que se busca, y está documentada como tal, no es
    un efecto colateral de otra cosa.
    \item \advertencia{Riesgo medio} \textbf{Compatibilidad de
    \texttt{response\_format strict} y visión pendiente de verificar} cuando el
    modelo pinneado es el mismo proveedor (p.~ej. \texttt{openai/gpt-4.1}), el
    riesgo es prácticamente nulo porque la petición llega sin traducir al
    endpoint real de OpenAI. El riesgo aparece si el fallback salta a un
    proveedor distinto (p.~ej. Anthropic), ya que cada proveedor tiene su
    propio mecanismo de salida estructurada y la traducción entre ambos debe
    probarse antes de confiar en ella en producción.
\end{itemize}

\section{Opción 3: LiteLLM (proxy autoalojado)}

\subsection{Qué es}
Proxy de código abierto, desplegado en infraestructura propia, que expone una
única API compatible con OpenAI por delante de más de 100 proveedores.

\subsection{Evaluación}

\begin{itemize}
    \item \ok{Sin intermediario que cobre margen.} Al ejecutarse en
    infraestructura propia, se paga directamente a cada proveedor (OpenAI,
    Anthropic, Azure...) su tarifa pública, sin comisión de un tercero.
    \item \ok{Máximo control y verificabilidad.} Al ser código abierto, se
    puede leer exactamente cómo traduce \texttt{response\_format} y las
    entradas de visión entre proveedores, en vez de confiar en el
    comportamiento de una caja negra.
    \item \ok{Fallback declarativo.} Cadenas de reintento entre modelos/
    proveedores configurables en un único archivo \texttt{config.yaml}, sin
    cambios en el código de \texttt{logic.py}/\texttt{imagenes.py} más allá de
    apuntar \texttt{OPENAI\_BASE\_URL} al proxy local.
    \item \riesgo{Coste operativo} Es una pieza de infraestructura más que hay
    que desplegar, actualizar y monitorizar (encaja con el stack actual de
    IIS/Windows Server/Python, pero añade mantenimiento propio que las otras
    dos opciones no requieren).
    \item \advertencia{Riesgo medio} Misma verificación pendiente que
    OpenRouter para la traducción de \texttt{response\_format strict} y visión
    al saltar entre familias de proveedor distintas.
\end{itemize}

\section{Comparativa}

\begin{table}[h]
\centering
\small
\begin{tabular}{p{3cm}p{3.3cm}p{3.3cm}p{3.3cm}}
\toprule
& \textbf{Azure OpenAI} & \textbf{OpenRouter} & \textbf{LiteLLM autoalojado} \\
\midrule
Compatibilidad \texttt{strict}/visión & Total (misma API) & Segura si se fija el mismo proveedor & Segura si se fija el mismo proveedor, verificable por ser código abierto \\
\addlinespace
Diversifica familia de modelo & No (mismo GPT) & Sí & Sí \\
\addlinespace
Fallback automático & No incluido de fábrica & Sí, nativo & Sí, declarativo \\
\addlinespace
Facturación & Pago por consumo, tarifa publicada & Pago por consumo + 5,5\% comisión publicada & Pago por consumo, sin comisión de terceros \\
\addlinespace
Esfuerzo de puesta en marcha & Mínimo (variables de entorno) & Bajo (cuenta + clave) & Medio (desplegar y mantener el proxy) \\
\addlinespace
Encaje con \texttt{OPENAI\_BASE\_URL} actual & Directo & Directo & Directo (apuntando al proxy local) \\
\bottomrule
\end{tabular}
\caption{Comparativa de las tres alternativas frente al objetivo de seguir
funcionando si OpenAI se cae.}
\end{table}

\section{Recomendación}

\begin{enumerate}
    \item \textbf{Corto plazo, mínimo esfuerzo:} activar \textbf{Azure OpenAI}
    como respaldo inmediato. Es el mismo modelo, la misma API, sin nada que
    probar ni que pueda romperse, y cubre el escenario más probable (caída del
    endpoint de \texttt{openai.com}).
    \item \textbf{Si además se quiere protección ante un problema específico de
    los modelos GPT} (no solo de la infraestructura de OpenAI), añadir
    \textbf{OpenRouter} como segunda capa, con un modelo de otra familia (p.~ej.
    Claude) configurado como fallback --- probando antes, con una factura de
    control, que la extracción con \texttt{json\_schema strict} funciona igual
    contra ese modelo concreto.
    \item \textbf{LiteLLM autoalojado} queda como opción a futuro si el volumen
    o el número de proveedores crece lo suficiente como para justificar
    mantener una pieza de infraestructura propia a cambio de eliminar cualquier
    comisión de terceros.
\end{enumerate}

No se recomienda Mammouth.ai para este objetivo: su modelo de cuotas por sesión
con degradación automática de modelo y su coste (superior al actual incluso en
el escenario más favorable) lo hacen menos adecuado que cualquiera de las tres
opciones anteriores, tal como se detalla en el informe previo
(\texttt{informe\_mammouth\_ai.tex}).

\end{document}
